\documentclass[11pt,a4paper]{article}
\usepackage[margin=2.0cm]{geometry}
\usepackage[T1]{fontenc}
\usepackage{lmodern}
\usepackage{amsmath}
\usepackage{booktabs}
\usepackage{tabularx}
\usepackage{xcolor}
\usepackage{listings}
\usepackage[hidelinks]{hyperref}

\emergencystretch=2em
\definecolor{codebg}{gray}{0.95}
\lstset{language=C++, basicstyle=\ttfamily\footnotesize, keywordstyle=\color{blue!70!black},
  commentstyle=\color{green!50!black}\itshape, backgroundcolor=\color{codebg},
  frame=single, rulecolor=\color{black!20}, breaklines=true, showstringspaces=false,
  tabsize=2, columns=flexible}
\newcommand{\code}[1]{\texttt{#1}}

\title{ATmega32A Module \\ \Large UART (Serial Communication)}
\author{Ali Sahafi}
\date{}

\begin{document}
\maketitle

\section{Theory}

UART (\emph{Universal Asynchronous Receiver and Transmitter}) sends data one
bit after the other. It uses one wire to send (\code{TXD}) and one wire to
receive (\code{RXD}), so it can send and receive at the same time (full
duplex). There is \emph{no clock wire}. For this reason, both sides must
use the same settings before they start:

\begin{itemize}
  \item the same \textbf{speed} (the \emph{baud rate}, in bits per second), and
  \item the same \textbf{frame}: how many data bits, parity or no parity, and
        how many stop bits.
\end{itemize}

\paragraph{The frame.} When nothing is sent, the line is HIGH (idle). Every
byte is sent as one frame:

\begin{center}
\begin{tabular}{llll}
\toprule
\textbf{Start bit} & \textbf{Data bits} & \textbf{Parity bit} & \textbf{Stop bit(s)} \\
\midrule
always 0 (LOW) & 8 bits, LSB first & optional & always 1 (HIGH), 1 or 2 bits \\
\bottomrule
\end{tabular}
\end{center}

The most common setting is \textbf{8N1}: 8 data bits, \textbf{N}o parity,
\textbf{1} stop bit. So one byte needs 10 bits on the wire.

\paragraph{Baud rate.} At 9600 baud, one bit takes $1/9600 \approx 104\,\mu$s.
The AVR makes the baud rate from the CPU clock with the register \code{UBRR}:
\[ \text{UBRR} = \frac{F_{\text{CPU}}}{16 \cdot \text{Baud}} - 1 \qquad\qquad
   \text{Baud} = \frac{F_{\text{CPU}}}{16 \cdot (\text{UBRR} + 1)} \]
\code{UBRR} must be a whole number, so the real baud rate is a little
different from the one you want. The error should be smaller than about 2\,\%.
At 8\,MHz:

\begin{center}
\begin{tabular}{rrrl}
\toprule
\textbf{Wanted baud} & \textbf{UBRR} & \textbf{Real baud} & \textbf{Error} \\
\midrule
4800   & 103 & 4808   & 0.2\,\% \\
9600   & 51  & 9615   & 0.2\,\% \quad (we use this) \\
19200  & 25  & 19231  & 0.2\,\% \\
38400  & 12  & 38462  & 0.2\,\% \\
57600  & 16 (double speed) & 58824  & 2.1\,\% \quad (too high) \\
115200 & 8 (double speed)  & 111111 & $-3.5$\,\% \quad (too high, does not work well) \\
\bottomrule
\end{tabular}
\end{center}

\paragraph{ASCII.} UART sends bytes (numbers 0--255). To send \emph{text}, every
character has a number: the ASCII code. For example \code{'A'} = 65,
\code{'a'} = 97, \code{'0'} = 48, \code{'1'} = 49, new line \code{'\textbackslash n'} = 10,
carriage return \code{'\textbackslash r'} = 13. A terminal program on the PC shows
each received byte as the character with that code. So if you send the
\emph{byte} 65, the terminal shows \code{A}. If you want the terminal to show
the \emph{number} 65, you must send the two characters \code{'6'} and \code{'5'}.

\section{How to use the driver}

The driver gives you one object: \code{UART}. You start it with one function call:

\begin{lstlisting}
UART.begin(9600);                        // 9600 baud, 8N1
UART.begin(9600, UART_PARITY_EVEN, 2);   // 9600 baud, 8 data bits, even parity, 2 stop bits
\end{lstlisting}

\subsection*{The options}

\begin{center}
\small
\begin{tabularx}{\textwidth}{@{}l l X@{}}
\toprule
\textbf{Argument} & \textbf{Value} & \textbf{Meaning} \\
\midrule
\code{baudRate} & \code{9600}, \code{19200}, \dots & Speed in bits per second. Write the number you want.
      The driver calculates \code{UBRR} for you. \\
\midrule
\code{parity}   & \code{UART\_PARITY\_NONE} (default) & No parity bit \\
                & \code{UART\_PARITY\_EVEN}           & Even parity: the number of 1s (data + parity) is even \\
                & \code{UART\_PARITY\_ODD}            & Odd parity: the number of 1s (data + parity) is odd \\
\midrule
\code{stopBits} & \code{1} (default), \code{2} & Number of stop bits \\
\bottomrule
\end{tabularx}
\end{center}

The driver always uses 8 data bits. The other side (for example the terminal
program on the PC) must use the \emph{same} settings.

\subsection*{Functions}

\begin{center}
\small
\begin{tabularx}{\textwidth}{@{}l X@{}}
\toprule
\textbf{Function} & \textbf{What it does} \\
\midrule
\code{UART.begin(baud, parity, stopBits)} & Set up and start the UART (see the options above) \\
\midrule
\code{UART.write(byte)} & Send one byte exactly as it is. \code{UART.write(65)} sends the byte 65, and the terminal shows \code{A} \\
\code{UART.print("text")} & Send text \\
\code{UART.print('A')} & Send one character \\
\code{UART.print(number)} & Send a number \emph{as text}: \code{UART.print(65)} sends the two characters \code{6} and \code{5}.
       Works for all integer types and for \code{float} (2 decimals; \code{UART.print(x, 4)} gives 4) \\
\code{UART.println(...)} & Same as \code{print()}, then a new line. \code{UART.println()} sends only the new line \\
\midrule
\code{UART.available()} & Returns \code{true} if a received byte is waiting. It does not wait \\
\code{UART.read()} & Waits until a byte has arrived, then returns it \\
\midrule
\code{UART.onReceive(function)} & Call \code{function(byte)} for every received byte (interrupt, see the interrupt lecture) \\
\bottomrule
\end{tabularx}
\end{center}

\paragraph{\code{write()} or \code{print()}?} Use \code{write()} to send raw
data (one byte, any value 0--255). Use \code{print()} and \code{println()} to
send text that a person reads in the terminal.

\paragraph{Waiting or not waiting.} \code{UART.read()} \emph{waits} until a byte
arrives: your program stops on this line until the user types something. If
the program must do other work at the same time, first ask
\code{UART.available()}, and only call \code{read()} when it returns \code{true}:
\begin{lstlisting}
if (UART.available()) {          // something has arrived?
  uint8_t data = UART.read();    // then read it (no waiting)
  ...
}
\end{lstlisting}

\paragraph{Interrupts (we come back to this in the interrupt lecture).}
\code{onReceive()} turns on the receive interrupt and calls \code{sei()} for
you. Your function then runs automatically for every received byte. Keep this
function short. If it changes a variable that \code{main()} also uses, declare
the variable as \code{volatile}. To turn the interrupt off again, pass
\code{nullptr}.

\subsection*{Driver functions and registers}

In the lecture slides we set the registers by hand. Each driver function does
these steps for you:

\begin{center}
\small
{\footnotesize\begin{tabular}{@{}ll@{}}
\toprule
\textbf{Driver function} & \textbf{Registers} \\
\midrule
\code{UART.begin(9600)} & \code{UBRRH = 0; UBRRL = 51;} \quad (UBRR = 8\,MHz / (16 $\cdot$ 9600) $-$ 1) \\
                        & \code{UCSRC = (1<{}<URSEL) | (1<{}<UCSZ1) | (1<{}<UCSZ0);} \quad (8N1) \\
                        & \code{UCSRB = (1<{}<RXEN) | (1<{}<TXEN);} \\
\code{UART.begin(9600, UART\_PARITY\_EVEN, 2)} & in \code{UCSRC} also: \code{UPM1} = 1 (even parity), \code{USBS} = 1 (2 stop bits) \\
\code{UART.write(x)} & \code{while (!(UCSRA \& (1<{}<UDRE)));  UDR = x;} \\
\code{UART.available()} & \code{UCSRA \& (1<{}<RXC)} \\
\code{UART.read()} & \code{while (!(UCSRA \& (1<{}<RXC)));  return UDR;} \\
\code{UART.onReceive(f)} & \code{UCSRB |= (1<{}<RXCIE); sei();} \\
\bottomrule
\end{tabular}}
\end{center}

\paragraph{Double speed.} If the normal \code{UBRR} gives an error of more than
1\,\%, the driver tries the double-speed mode (\code{U2X} in \code{UCSRA},
$\text{UBRR} = F_{\text{CPU}} / (8 \cdot \text{Baud}) - 1$) and uses it if the
error is smaller. At 8\,MHz this only happens for 57600 and 115200 baud.

\newpage
\section{Example: echo with ASCII code}

\paragraph{Board hardware.} The UART pins \code{PD0} (\code{RXD}) and \code{PD1}
(\code{TXD}) are connected to the \textbf{FT232} chip on the board. This chip
is a USB-to-UART converter. When you connect the board's USB cable to the PC,
the PC gets a new serial port (a \emph{virtual COM port}):
\begin{itemize}
  \item \textbf{Windows:} \code{COM3}, \code{COM4}, \dots\ (see \emph{Device Manager $\rightarrow$ Ports}).
  \item \textbf{macOS:} \code{/dev/cu.usbserial-...}
  \item \textbf{Linux:} \code{/dev/ttyUSB0}
\end{itemize}
You also need the LEDs (D0--D7 on \code{PORTB}, active-low: \code{LOW} = ON)
for Task 1.

\paragraph{Serial terminal.} To see the data on the PC, you need a terminal
program, for example the \emph{Serial Monitor} extension in VS Code, CoolTerm
or PuTTY. Choose the serial port of the board and these settings:
\textbf{9600 baud, 8 data bits, no parity, 1 stop bit (8N1)}, no flow control.

\paragraph{Build and upload.} \code{make uart}. Then open the terminal and type
a character, for example \code{A}. The board answers
\code{You typed: A  (ASCII 65)}.

\lstinputlisting{example.cpp}

\newpage
\section{Laboratory Tasks}
\setlength{\parindent}{0pt}\setlength{\parskip}{5pt}

For each task, make a new \code{.cpp} file in the main project folder (for
example \code{Task1.cpp}). Start the file with \code{\#define F\_CPU 8000000UL}.
Then include the drivers you need: \code{drivers/uart/uart.hpp},
\code{drivers/gpio/gpio.hpp}, \code{drivers/timer/timer.hpp}. Upload it with
\code{make SRC=Task1.cpp}. Use 9600 baud (8N1) in all tasks.

\subsection*{Task 1: Send and Receive Bytes (20 min)}
\begin{itemize}
  \item \textbf{Task 1.1 (Send):} Send the numbers 0 to 255 to the PC, one number
        every second. Use \code{UART.write(number)} and \code{\_delay\_ms(1000)}.
        Look at the data in the terminal. What do you see? Then switch the
        terminal to show the received bytes in HEX (for example in CoolTerm).
        Now what do you see? (The ASCII part of the lecture explains the difference.)
  \item \textbf{Task 1.2 (Receive):} Wait for a byte from the PC with
        \code{UART.read()} and show it on the 8 LEDs on \code{PORTB}. The LEDs are
        active-low, so write the inverted value:
        \mbox{\code{GPIO.write(PORTB, \textasciitilde data)}}. Type the character \code{A}
        in the terminal. Which LEDs are on? Compare with the ASCII code of \code{A}
        in binary. Try \code{a} and \code{1} too.
\end{itemize}

\subsection*{Task 2: Send Your Name (ASCII Text)}
Send your name to the PC, one time every second. Do it in two steps:
\begin{enumerate}
  \item First send only one character every second, for example
        \code{UART.print('A')}.
  \item Then send your whole name and a new line with
        \code{UART.println("Your Name")}.
\end{enumerate}
\code{println()} adds the characters \code{'\textbackslash r'} and
\code{'\textbackslash n'}, so every name starts on a new line in the terminal.

\subsection*{Assignment: Send and Receive at the Same Time}
Write one program that does two things at the same time:
\begin{itemize}
  \item It sends the numbers 0 to 255 to the PC as text with
        \code{UART.println(number)}, one number every second.
  \item It shows every byte that it receives from the PC on the LEDs on
        \code{PORTB} (active-low).
\end{itemize}
Rules:
\begin{itemize}
  \item Use \code{Timer0} \emph{without interrupts} for the 1 second, for example
        with a 10\,ms tick in \code{TIMER\_CTC} mode and
        \code{Timer0.compareMatched()} (see the timer handout). Count 100 ticks
        for one second. Do \emph{not} use \code{\_delay\_ms()}.
  \item Check for received data all the time with \code{UART.available()}.
        Do \emph{not} call \code{UART.read()} without asking
        \code{UART.available()} first, or the program stops and waits.
\end{itemize}

The main loop then looks like this:
\begin{lstlisting}
while (true) {
  if (Timer0.compareMatched()) {   // 10 ms have passed
    ...                            // count ticks; after 100 ticks send the next number
  }
  if (UART.available()) {          // a byte has arrived
    ...                            // read it and show it on the LEDs
  }
}
\end{lstlisting}

\end{document}
